% Related Work -- fractional/linear programming subsection.
% Revision of the "Franctional programing" paragraph. New citation keys are listed in
% related-fractional.bib; keep your own keys if they already exist.

\paragraph{Fractional and linear programming for distributional constraints.}
Stage 1 maximizes a ratio of two linear functions of binary decision variables, placing it in
the class of \emph{fractional $0$--$1$ programs} (also called hyperbolic $0$--$1$ programs),
surveyed by Borrero et al.~\cite{borrero2017fractional}. For continuous variables, Charnes and
Cooper~\cite{charnes1962programming} showed that a linear-fractional program can be linearized
by a change of variables and solved as a single linear program, but this transformation does
not survive the integrality restriction. We therefore use the parametric method of
Dinkelbach~\cite{dinkelbach1967nonlinear}, which replaces the ratio with a sequence of linear
objectives of the form $N(S) - \lambda\,n(S)$ and updates $\lambda$ to the ratio achieved by
the incumbent solution. Each iteration is thus an ordinary linear maximization over the same
feasible set---in our case a single top-$\alpha k$ selection---and the sequence converges to
the exact optimum. The rapid convergence we observe empirically (Section~\ref{sec:dinkelbach})
is consistent with what is known for the combinatorial case: Megiddo's parametric
search~\cite{megiddo1979combinatorial} and Radzik's analysis of Newton's
method~\cite{radzik1992newton}, of which Dinkelbach's method is an instance, establish
strongly polynomial bounds on the number of iterations required for fractional combinatorial
optimization.

Constraints of this ratio form recur throughout the fairness literature, since most group
fairness criteria compare selection or error rates \emph{between} groups. Celis et
al.~\cite{DBLP:conf/fat/CelisHKV19} give a meta-algorithm for classification subject to any
fairness metric expressible as a linear-fractional constraint, which covers statistical parity,
false discovery rate, and related group measures, and Chen and
Hooker~\cite{chen2021guide} survey the wider space of fairness and equity criteria in
optimization models, including proportional and ratio-based formulations. Our setting differs
in where the ratio arises. In these works the ratio is evaluated over a fixed dataset or a
fixed ground set of decisions, whereas in DUTS each selected table contributes both matching
tuples and total tuples, so the numerator and the denominator of the achieved proportion are
both determined by the selection itself. The distributional requirement is therefore genuinely
fractional rather than a linear constraint in disguise, which is what makes the fractional
programming machinery necessary rather than merely convenient. To the best of our knowledge,
this is the first application of fractional programming to table search over data lakes.
